\documentclass[11pt]{article}

\usepackage[utf8]{inputenc}
\usepackage[T1]{fontenc}
\usepackage{charter}
\usepackage{amsmath,amssymb}
\usepackage{booktabs}
\usepackage{array}
\usepackage{graphicx}
\usepackage{microtype}
\usepackage{setspace}
\usepackage{cite}
\usepackage[hidelinks]{hyperref}
\usepackage[a4paper,margin=2.5cm]{geometry}
\usepackage[font=small,labelfont=bf]{caption}
% Auto-generated by analysis/build_public_core_macros.py; do not edit.
% Source: results/public_core_evidence.json sha256=c41aff9fa9deeaac8a77d43e95b4029ca1e9ad3565f465c35e73e15741426caf
\newcommand{\DffRows}{12,651}
\newcommand{\DffTargets}{44}
\newcommand{\DsRows}{260,060}
\newcommand{\DsTargets}{58}
\newcommand{\DffRawPR}{1.834}
\newcommand{\DffResidualPR}{9.303}
\newcommand{\DsRawPR}{2.266}
\newcommand{\DsResidualPR}{18.206}
\newcommand{\DffRawMeanRtwo}{0.535}
\newcommand{\DffResidualMeanRtwo}{0.087}
\newcommand{\DsRawMeanRtwo}{0.431}
\newcommand{\DsResidualMeanRtwo}{0.038}
\newcommand{\DffAxisCosine}{0.970}
\newcommand{\DsAxisCosine}{0.969}
\newcommand{\DffCovRawPR}{1.915}
\newcommand{\DffCovResidualPR}{6.911}
\newcommand{\DsCovRawPR}{2.474}
\newcommand{\DsCovResidualPR}{8.235}
\newcommand{\DffNullMedianLow}{42.400}
\newcommand{\DffNullMedianHigh}{42.628}
\newcommand{\DsNullMedianLow}{56.481}
\newcommand{\DsNullMedianHigh}{56.612}
\newcommand{\DffMWConditionedNullPR}{34.875}
\newcommand{\DsMWConditionedNullPR}{47.039}
\newcommand{\DffMWConditionedExplainedShare}{0.238}
\newcommand{\DsMWConditionedExplainedShare}{0.251}
\newcommand{\DffMWConditionedMapRho}{0.924}
\newcommand{\DsMWConditionedMapRho}{0.718}
\newcommand{\DffMWConditionedSignAgreement}{0.834}
\newcommand{\DsMWConditionedSignAgreement}{0.727}
\newcommand{\DffMWQuartileRho}{0.205}
\newcommand{\DsMWQuartileRho}{0.351}
\newcommand{\DffMWMatchedControlRho}{0.992}
\newcommand{\DsMWMatchedControlRho}{0.994}
\newcommand{\DffMWContinuumRho}{0.945}
\newcommand{\DsMWContinuumRho}{0.988}
\newcommand{\DffSizeAxisRhoRange}{0.200--0.207}
\newcommand{\DsSizeAxisRhoRange}{0.322--0.361}
\newcommand{\DffOtherAxisRhoRange}{0.703--0.887}
\newcommand{\DsOtherAxisRhoRange}{0.787--0.831}
\newcommand{\DffProbeTwoHundred}{0.940}
\newcommand{\DffProbeFiveHundred}{0.973}
\newcommand{\DsProbeTwoHundred}{0.920}
\newcommand{\DsProbeFiveHundred}{0.965}
\newcommand{\DffProbeTwoHundredSign}{0.896}
\newcommand{\DffProbeFiveHundredSign}{0.931}
\newcommand{\DsProbeTwoHundredSign}{0.877}
\newcommand{\DsProbeFiveHundredSign}{0.920}
\newcommand{\DffScaffoldProbeTwoHundred}{0.925}
\newcommand{\DffScaffoldProbeFiveHundred}{0.958}
\newcommand{\DsScaffoldProbeTwoHundred}{0.914}
\newcommand{\DsScaffoldProbeFiveHundred}{0.958}
\newcommand{\DffRandomHoldoutTwoHundred}{0.935}
\newcommand{\DffRandomHoldoutFiveHundred}{0.968}
\newcommand{\DsRandomHoldoutTwoHundred}{0.917}
\newcommand{\DsRandomHoldoutFiveHundred}{0.960}
\newcommand{\DffSelectorRawRtwo}{0.900}
\newcommand{\DsSelectorRawRtwo}{0.698}
\newcommand{\DffSelectorResidualRtwo}{0.278}
\newcommand{\DsSelectorResidualRtwo}{0.219}
\newcommand{\DffSelectorDirectResidualRtwo}{0.278}
\newcommand{\DsSelectorDirectResidualRtwo}{0.219}
\newcommand{\DffSelectorResidualPCOneRtwo}{0.180}
\newcommand{\DsSelectorResidualPCOneRtwo}{0.065}
\newcommand{\DffSelectorDirectResidualDelta}{0.000015}
\newcommand{\DsSelectorDirectResidualDelta}{-0.000001}
\newcommand{\DffSelectorCommonResidualDelta}{0.036}
\newcommand{\DsSelectorCommonResidualDelta}{0.058}
\newcommand{\DffPanelTwoHundredDelta}{0.0153}
\newcommand{\DffPanelFiveHundredDelta}{0.0083}
\newcommand{\DsPanelTwoHundredDelta}{0.0113}
\newcommand{\DsPanelFiveHundredDelta}{0.0056}
\newcommand{\ExperimentalCrossRho}{0.848}
\newcommand{\ExperimentalCrossRawRho}{0.582}
\newcommand{\ExperimentalCrossQAP}{0.0001}
\newcommand{\DavisSplitRho}{0.757}
\newcommand{\PkistwoSplitRho}{0.915}
\newcommand{\DavisSameSupportLigands}{56}
\newcommand{\PkistwoSameSupportLigands}{154}
\newcommand{\DavisSameSupportRho}{0.307}
\newcommand{\PkistwoSameSupportRho}{0.326}
\newcommand{\DavisBroadSupportRho}{0.300}
\newcommand{\PkistwoBroadSupportRho}{0.329}
\newcommand{\HotspotLigands}{176}
\newcommand{\HotspotTargets}{21}
\newcommand{\HotspotSplitRho}{0.794}
\newcommand{\HotspotSplitRhoLow}{0.690}
\newcommand{\HotspotSplitRhoHigh}{0.854}
\newcommand{\DavisHotspotRawRho}{0.690}
\newcommand{\DavisHotspotResidualRho}{0.765}
\newcommand{\PkistwoHotspotRawRho}{0.526}
\newcommand{\PkistwoHotspotResidualRho}{0.751}
\newcommand{\PanelTransferPositiveCells}{14}
\newcommand{\PanelTransferTotalCells}{15}
\newcommand{\PanelTransferConservativeMean}{0.0607}
\newcommand{\PanelTransferJointAUC}{0.0616}
\newcommand{\PanelTransferJointP}{0.0069}
\newcommand{\PanelRankPositiveCells}{14}
\newcommand{\PanelRankConservativeMean}{0.0384}
\newcommand{\PanelRankP}{0.0056}
\newcommand{\PanelTargetLooPositive}{21}
\newcommand{\PanelTargetLooTotal}{21}
\newcommand{\PanelTargetLooAucLow}{0.0294}
\newcommand{\PanelTargetLooAucHigh}{0.0622}
\newcommand{\PanelTargetLooCellWinsLow}{11}
\newcommand{\PanelTargetLooCellWinsHigh}{15}
\newcommand{\HotspotClusterMedian}{0.0628}
\newcommand{\HotspotClusterLow}{0.0319}
\newcommand{\HotspotClusterHigh}{0.1063}
\newcommand{\HotspotClusterPositiveFraction}{1.000}
\newcommand{\DffMedianImputeResidualPR}{9.528}
\newcommand{\DffMissingAwareResidualPR}{9.663}
\newcommand{\DffCompleteCaseResidualPR}{13.970}
\newcommand{\DffEqualNResidualPRLow}{9.128}
\newcommand{\DffEqualNResidualPRHigh}{9.445}
\newcommand{\SharedExactLigands}{582}
\newcommand{\SharedConnectivityBlocks}{804}
\newcommand{\SharedTargetIdentities}{9}
\newcommand{\SharedReceptorStructures}{3}


\hypersetup{
  pdftitle={Supplementary information: Two-way centering reveals distributed target-relative structure and limits score-surface compression in Vina docking matrices},
  pdfauthor={Adeliya Leleytner, Victor Safronov, Maxim Fedorov}
}

\renewcommand{\thetable}{S\arabic{table}}
\renewcommand{\thefigure}{S\arabic{figure}}
\renewcommand{\arraystretch}{1.12}
\setlength{\emergencystretch}{3em}
\doublespacing

\title{Supplementary information\\[4pt]
\large Two-way centering reveals distributed target-relative structure and
limits score-surface compression in Vina docking matrices}
\author{Adeliya Leleytner, Victor Safronov, Maxim Fedorov}
\date{}

\begin{document}
\maketitle

\noindent\textbf{Scope.}
This Supplement is restricted to the public sources used by the public-core
manuscript: Docking-44, DOCKSTRING, DAVIS-Complete, PKIS2 and the
checksum-fetched Anastassiadis HotSpot workbook. It contains only the
matrix-spectrum, chemical-support, pilot-recovery, score-reconstruction,
experimental-panel and reusable-audit analyses assembled in
\path{results/public_core_evidence.json}. Ancillary public identity and target
mapping files supply identifiers for the overlap audit but no additional score
or assay rows. Historical or source-restricted analysis branches are outside
the report layer and are not summarized here.

\section*{S1. Data, estimands and cross-panel overlap}

\subsection*{S1.1 Public row-level sources}

Table~\ref{tab:s1data} states the support actually used by this Supplement.
Docking-44 was taken from \path{data/frozen/df_final_v4.csv.gz}
\cite{nikitin2025}. Its \DffRows{} rows contain \DffTargets{} Vina-family
target columns. The primary preprocessing clips positive scores at zero (none
occurred) and replaces a missing score with the mean of its target. DOCKSTRING
was taken from \path{data/frozen/dockstring-dataset.tsv.gz}
\cite{dockstring2022}. Of 260,155 source rows, 95 contained at least one of 339
missing target scores and were removed, leaving \DsRows{} complete rows by
\DsTargets{} targets. The released analysis contract clips its 5,393 positive
scores to zero; this is 0.0358\% of cells in the complete matrix.

DAVIS-Complete is redistributed as
\path{data/frozen/davis_complete.tab.gz} \cite{davis2011,wu2025davis}. The
source audit recovers the complete 72-compound by 442-target grid (31,824
records). PKIS2 is redistributed as \path{data/frozen/pkis2_s4.xlsx}
\cite{drewry2017}; it contains 645 populated compound rows and 406 assay
columns. The public-core experimental analysis restricts both panels to the
same 21 kinases present in DOCKSTRING. This gives 13,545 observed PKIS2 cells
with no missing values. Five DAVIS profiles and one PKIS2 profile are constant
over these 21 targets; the remaining 67 and 644 profiles define the
informative cross-panel maps.

Anastassiadis Supplementary Table~3 is downloaded from the publisher and
accepted only at the frozen SHA-256 \cite{anastassiadis2011}. It is not
redistributed. The primary 21-target block contains 176 complete compounds
after two declared incomplete profiles are removed; CDK2 is the mean of its
cyclin-A and cyclin-E measurements.

\begin{table}[htbp]
\centering
\caption{Public inputs and fixed analysis supports. Licences and full
SHA-256 digests are recorded in the public-core source manifest.}
\label{tab:s1data}
\small
\begin{tabular}{p{0.18\linewidth}p{0.30\linewidth}p{0.43\linewidth}}
\toprule
Source & Redistributed file & Public-core support \\
\midrule
Docking-44 & \path{data/frozen/df_final_v4.csv.gz} &
12,651 ligands $\times$ 44 targets; 8,159 missing cells target-mean imputed \\
DOCKSTRING & \path{data/frozen/dockstring-dataset.tsv.gz} &
260,060 complete ligands $\times$ 58 targets; 95 incomplete source rows removed \\
DAVIS-Complete & \path{data/frozen/davis_complete.tab.gz} &
72 ligands $\times$ 442 targets in the source audit; 21 kinases in the public-core map \\
PKIS2 & \path{data/frozen/pkis2_s4.xlsx} &
645 compounds; 21 complete kinase columns in the public-core map \\
HotSpot & Publisher Supplementary Table~3, checksum-fetched only &
176 complete compounds; 21 kinase columns in the public-core map \\
\bottomrule
\end{tabular}
\end{table}

Input identities, byte counts, licences and SHA-256 digests are machine-readable
in \path{results/public_core_source_manifest.csv}; the digest of that manifest
is stored in \path{results/public_core_source_manifest.sha256}. The public-core
builder fails if an opened source is undeclared or if its digest has changed.
The overlap audit additionally reads the frozen identifier contract
\path{data/frozen/dockstring_identity_contract_2026-08-03.csv.gz}, the public
family labels in \path{data/target_families.csv}, and a curated receptor mapping
whose rows and source digest are exposed in
\path{results/cross_panel_overlap_audit/summary.json}. These files define
identities and labels only; they add no ligand--target measurements.

\subsection*{S1.2 Map and dimension estimands}

For a ligand-by-target score or response matrix $X$, the two-way residual is
\begin{equation}
\Gamma_{ij}=X_{ij}-\bar X_{i\cdot}-\bar X_{\cdot j}+\bar X_{\cdot\cdot}.
\label{eq:s1residual}
\end{equation}
The target map is the Pearson correlation matrix of the specified columns, and
its geometry is the vector of unique off-diagonal correlations. Correlation
spectra standardize each target to unit variance; covariance spectra retain
target-specific variances. For a $P$-target correlation matrix $C$ with
eigenvalues $\lambda_k$, the participation-ratio dimension is
\begin{equation}
r_{\mathrm{PR}}
=\frac{(\sum_k\lambda_k)^2}{\sum_k\lambda_k^2}
=\frac{P}{1+(P-1)\overline{r_{ij}^{\,2}}}.
\label{eq:s1pr}
\end{equation}
Thus PR is an exact re-expression of mean squared target correlation, not an
algebraic rank or an independently estimated molecular dimension
\cite{mazzucato2016,roy2007}. Equation~\ref{eq:s1residual} also imposes one zero
spectral mode because every residual row sums to zero.

\subsection*{S1.3 Chemical and target overlap of the two Vina matrices}

The two Vina matrices were generated independently but are not disjoint
statistical samples. Recomputed full Standard InChIKeys identify 582 shared
ligands. Collapsing stereochemical and protonation suffixes gives 804 shared
connectivity blocks, containing 804 Docking-44 rows and 1,014 DOCKSTRING rows.
There are also 1,340 shared nonempty Bemis--Murcko scaffolds, containing 7,529
Docking-44 and 25,614 DOCKSTRING rows \cite{murcko1996}. The curated target map
contains nine matched receptor identities, although only seven Docking-44
structures are human. The 2vt4 and 3ln1 structures are turkey ADRB1 and mouse
PTGS2 orthologues; both are among the three mapped pairs that use the same
receptor structure in DOCKSTRING. The complete overlap record, including all
target mappings, is
\path{results/cross_panel_overlap_audit/summary.json}. Agreement across these
matrices is therefore replication across independently sourced score matrices,
not evidence from formally independent chemical or receptor samples.

\subsection*{S1.4 Interpretation of ranges and probabilities}

No interval in this Supplement is presented as a confidence interval for a
population of proteins or all chemical space. Central 95\% \emph{simulation
ranges} are empirical 2.5th--97.5th percentiles of null simulations. Central
95\% \emph{repeated-sample}, \emph{repeated-support} or
\emph{repeated-split ranges} are the same percentiles over the stated random
sampling operation. Percentile \emph{bootstrap perturbation ranges} condition
on the fixed target panel and observed assay systems. The one QAP probability
reported below is a two-sided Monte Carlo target-label probability. These
objects quantify different sources of variation and are not interchangeable.

\section*{S2. Spectral, covariance, missing-data and null sensitivities}

\subsection*{S2.1 Correlation and covariance spectra}

The full-support summaries are given in Table~\ref{tab:s2spectra}. In both
matrices the raw first component is close to the uniform target direction:
loading-vector cosine similarities are 0.970 for Docking-44 and 0.969 for
DOCKSTRING. The corresponding component scores correlate 0.998 and 0.997 with
the mean column-standardized ligand score. Two-way centring increases PR from
1.834 to 9.303 and from 2.266 to 18.206, while mean squared target correlation
falls from 0.535 to 0.087 and from 0.431 to 0.038.

\begin{table}[htbp]
\centering
\caption{Full-support correlation-spectrum summaries. The numbers are fixed
matrix summaries, not estimates for a target superpopulation.}
\label{tab:s2spectra}
\small
\resizebox{\textwidth}{!}{%
\begin{tabular}{lrrrrrr}
\toprule
Dataset & Raw PC1 & Raw PR & Raw $\overline{r^2}$ & Residual PR &
Residual $\overline{r^2}$ & PCs for 90\%, raw/residual \\
\midrule
Docking-44 & 0.733 & 1.834 & 0.535 & 9.303 & 0.087 & 8/29 \\
DOCKSTRING-58 & 0.659 & 2.266 & 0.431 & 18.206 & 0.038 & 22/43 \\
\bottomrule
\end{tabular}
}
\end{table}

Unit-variance scaling changes the magnitude but not the direction of this
contrast. On unscaled covariance surfaces, Docking-44 PR rises from 1.915 to
6.911 and DOCKSTRING PR from 2.474 to 8.235 after two-way centring. Retaining
the 5,393 positive DOCKSTRING scores instead of clipping them changes raw PR
from 2.266 to 2.275 and residual PR from 18.206 to 18.362. The complete scalar
records, including entropy dimension, PC1 fraction and the exact PR identity,
are in the \texttt{large\_vina\_panels} block of
\path{results/public_core_evidence.json}.

\subsection*{S2.2 Docking-44 missing scores alter support under complete-case deletion}

Docking-44 contains 8,159 missing scores among 556,644 cells (1.4657\%), spread
over 3,354 ligands. Target 4mqs accounts for 3,352 missing cells. Target-mean
imputation on all 12,651 rows gives residual PR 9.3029; target-median imputation
gives 9.5279. Fitting additive row and target effects only to the 548,485
observed cells and completing an unobserved residual with its fitted value of
zero gives 9.6630. This last construction is an additive-model sensitivity, not
recovery of an unknown raw docking score.

Complete-case deletion retains 9,297 ligands and gives residual PR 13.9702.
Restricting the primary mean-imputed surface to exactly those rows produces an
identical score matrix, confirming that no imputed value survives the
complete-case comparison. Across 1,000 equally sized random supports drawn from
the full primary surface, residual PR has median 9.2940 and central 95\%
repeated-support range 9.1281--9.4452. The complete-case value lies above every
random-support value.

The selected rows are chemically atypical: complete and excluded-incomplete
rows have median molecular weights of 249.332 and 398.552~g/mol and median
heavy-atom counts of 17 and 28. The molecular-weight median shift is
149.220~g/mol (Cliff's $\delta=0.893$), and the heavy-atom median shift is 11
($\delta=0.896$). These are descriptive associations with row inclusion; no
missing-at-random mechanism or causal explanation is assumed. Replicate-level
values and chemical summaries are in
\path{results/public_docking44_missing_data_sensitivity/equal_n_random_controls.csv},
\path{results/public_docking44_missing_data_sensitivity/spectral_sensitivity.csv},
\path{results/public_docking44_missing_data_sensitivity/chemistry_group_summary.csv}
and \path{results/public_docking44_missing_data_sensitivity/chemistry_effect_sizes.csv}.

\subsection*{S2.3 Transformation-matched residual nulls}

Four null families were evaluated on one deterministic 12,000-ligand support
per dataset, with 500 simulations per family. The additive Gaussian null keeps
fitted ligand and target effects and target-specific residual variances. The
empirical-column null independently permutes each observed residual target
column before repeating two-way centring and scaling. The row-norm null samples
an independent direction in the target-zero-sum subspace for each ligand and
rescales it to that ligand's observed residual Euclidean norm. It therefore
preserves the fixed support, target count, zero row sums and every row norm, but
not target marginals or cross-target alignment.

The fourth family independently permuted processed raw target scores within
empirical molecular-weight bins and then repeated row centring. At the
permutation step it preserves each target's exact empirical score marginal
within each bin but not ligand correspondence across targets. The ten-bin
design was primary; 1, 5 and 20 bins and random partitions with the same bin
sizes were sensitivities.

\begin{table}[htbp]
\centering
\caption{Residual PR on the common 12,000-ligand support and under four
matched nulls. Brackets are central 95\% simulation ranges over 500 null
replicates, not confidence intervals.}
\label{tab:s2nulls}
\small
\begin{tabular}{p{0.20\linewidth}p{0.14\linewidth}p{0.27\linewidth}p{0.27\linewidth}}
\toprule
Dataset & Observed & Null family & Null median [simulation range] \\
\midrule
Docking-44 & 9.295 & Additive Gaussian & 42.400 [42.345, 42.451] \\
 & & Residual-column permutation & 42.403 [42.350, 42.451] \\
 & & Row-norm random direction & 42.628 [42.597, 42.663] \\
 & & MW-conditioned raw-score permutation (10 bins) & 34.875 [34.790, 34.972] \\
\addlinespace
DOCKSTRING-58 & 18.158 & Additive Gaussian & 56.482 [56.445, 56.515] \\
 & & Residual-column permutation & 56.481 [56.438, 56.516] \\
 & & Row-norm random direction & 56.612 [56.583, 56.640] \\
 & & MW-conditioned raw-score permutation (10 bins) & 47.039 [46.942, 47.131] \\
\bottomrule
\end{tabular}
\end{table}

Every simulated residual PR exceeded its corresponding observed value. For
each family the lower-tail Monte Carlo probability is therefore
$(0+1)/(500+1)=0.0020$, the resolution floor of that simulation design. The
first three nulls reject three mechanical explanations for the residual
spectral concentration, and the fourth tests coarse MW-conditioned marginals;
none establishes docking accuracy, biological
correctness or transport beyond the two fixed panels. Seeds, support-index
digests, minima, maxima and all distribution summaries are in
\path{results/public_residual_null_audit/summary.json}; all conditional-null
replicates and exact bin contracts are in
\path{results/public_residual_null_audit/mw_conditional_null_repeats.csv} and
\path{results/public_residual_null_audit/mw_conditional_bin_contract.csv}.

Across 5, 10 and 20 bins, the share of the one-bin-to-observed PR gap explained
by MW conditioning was 0.191, 0.238 and 0.257 for Docking-44 and 0.228, 0.251
and 0.260 for DOCKSTRING-58. The corresponding ten-bin size-matched random
partitions remained at PR 42.844 and 56.726. Although ten-bin MW conditioning
left PR far above observation, its median Spearman agreement with the observed
residual edge map was \DffMWConditionedMapRho{} and
\DsMWConditionedMapRho{}, with sign agreement
\DffMWConditionedSignAgreement{} and \DsMWConditionedSignAgreement{}. Thus the
coarse size gradients recover substantial map topology while accounting for
only about one quarter of its excess concentration.

\section*{S3. Chemical-domain controls}

\subsection*{S3.1 Quartile contrast and matched chemical-group controls}

Molecular weight was selected after preliminary inspection and is treated as a
post-hoc domain coordinate. The primary Docking-44 contrast uses all 12,651
rows; its lower and upper threshold-defined quartiles contain 3,163 rows each,
with median molecular weights 178.054 and 413.602~g/mol. The primary
DOCKSTRING contrast uses a score-blind 15,000-row support drawn with seed 71;
its two bands contain 3,751 and 3,750 rows, with medians 286.310 and
498.573~g/mol.

Docking-44 uses 8,150 frozen source cluster labels from the
\texttt{Butina\_clusters} column; the public producer treats these labels as
opaque and does not regenerate or assume their fingerprint recipe. On the
primary DOCKSTRING support, cyclic molecules are grouped by Bemis--Murcko
scaffold and each acyclic molecule is a source-row singleton, yielding 11,905
groups \cite{butina1999,murcko1996}. The global matched control assigns whole
groups to disjoint halves within ten rank strata of group-median molecular
weight while balancing ligand counts. The within-band controls repeat the same
whole-group split after restricting to one molecular-weight band.

\begin{table}[htbp]
\centering
\caption{Residual-map agreement across molecular-weight domains and controls.
The observed low/high value is one fixed contrast. Control brackets are central
95\% repeated-split ranges across 200 chemical-group assignments, not
confidence intervals.}
\label{tab:s3controls}
\small
\resizebox{\textwidth}{!}{%
\begin{tabular}{p{0.18\linewidth}rrrr}
\toprule
Dataset & Observed low/high $\rho$ & Sign flips & Global MW-matched control &
Within low / within high \\
\midrule
Docking-44 & 0.205 & 0.425 &
0.992 [0.985, 0.995] &
0.980 [0.974, 0.984] / 0.979 [0.966, 0.986] \\
DOCKSTRING-58 & 0.351 & 0.367 &
0.994 [0.993, 0.995] &
0.979 [0.973, 0.982] / 0.966 [0.958, 0.971] \\
\bottomrule
\end{tabular}
}
\end{table}

At the fixed quartile definition, the observed agreement is below the 2.5th
percentile of the global and both within-band whole-group controls in each
matrix. The controls separate the low/high contrast from finite-sample
instability and generic chemical disjointness within a fixed molecular-weight
band. They do not isolate molecular weight from correlated chemical properties.
Raw-map results lead to the same boundary: observed low/high agreements are
0.227 and 0.409, compared with global matched-control medians 0.996 and 0.998.
The complete summaries and replicate rows are in
\path{results/public_chemical_domain_controls/observed_low_high_mw_contrasts.csv},
\path{results/public_chemical_domain_controls/mw_matched_group_disjoint_summary.csv}
and
\path{results/public_chemical_domain_controls/within_band_reproducibility_summary.csv}.

\subsection*{S3.2 Threshold, continuum and support-seed sensitivities}

A separate exact-rank sensitivity used lower and upper tail fractions of 0.10,
0.15, 0.20, 0.25, 0.30, 0.35 and 0.40. Residual-map agreement ranges from
0.201 to 0.259 in Docking-44 and from 0.174 to 0.528 in DOCKSTRING. At every
fraction, the observed value is below the 2.5th percentile of 100 equal-$N$
random-disjoint controls. Their mean agreements range from 0.978 to 0.995 and
from 0.976 to 0.994, respectively. These are threshold-sensitivity and
repeated-support ranges, not confidence intervals.

Ten non-overlapping equal-population molecular-weight bins provide a continuous
check. Across their 45 pairs, Spearman correlation between separation of bin
median molecular weights and residual-map dissimilarity is 0.945 for
Docking-44 and 0.988 for DOCKSTRING. The 45 observations reuse ten maps, so the
trend is descriptive and no independent-pair significance calculation is
reported. Exact threshold and bin-pair records are in
\path{results/revision_map_decision_sensitivities/mw_threshold_sweep.csv} and
\path{results/revision_map_decision_sensitivities/mw_decile_pair_distances.csv};
compact summaries are in
\path{results/revision_map_decision_sensitivities/mw_threshold_sweep_summary.csv}
and
\path{results/revision_map_decision_sensitivities/mw_decile_continuous_trends.csv}.

Finally, the DOCKSTRING quartile contrast was repeated on six independently
seeded, score-blind 15,000-ligand supports. Residual-map agreement ranges from
0.325 to 0.361 (median 0.346), while raw-map agreement ranges from 0.377 to
0.409 (median 0.398). This is a range across six specified support seeds, not
an uncertainty interval. Support-index SHA-256 fingerprints and per-seed
values are in
\path{results/public_chemical_domain_controls/dockstring_support_seed_contrasts.csv}.

\subsection*{S3.3 Descriptor-domain specificity}

To test whether the molecular-weight result generalized indiscriminately to
other coarse chemical coordinates, the same lower/upper 25\% threshold family
was applied post hoc to seven RDKit descriptors. Heavy-atom count, molecular
weight and Labute ASA were designated size-related; TPSA, cLogP, rotatable-bond
count and ring count formed the other-coarse family. Inclusive empirical
quartile thresholds retain boundary ties, so discrete descriptors need not
produce equal-sized tails. Docking-44 again uses all 12,651 rows and
DOCKSTRING uses the fixed seed-71 support of 15,000 rows.

Across the three size-related axes, residual-map agreement spans
0.200--0.207 (median 0.205) in Docking-44 and 0.322--0.361 (median 0.351) in
DOCKSTRING. Across the four other coarse axes, the corresponding spans are
0.703--0.887 (median 0.771) and 0.787--0.831 (median 0.812). These are
descriptive ranges across the named, correlated descriptors, not confidence
intervals or evidence that any descriptor is causal. The three size-related
coordinates are themselves strongly correlated and represent one coherent
family, not three independent replications.

For every descriptor, 200 random-row-disjoint controls used exactly the
observed low- and high-tail sizes. Every observed contrast fell below the 2.5th
percentile of its matched control (Table~\ref{tab:s3descriptors}). Thus the
result is graded rather than a size-versus-null dichotomy: the strongest changes
follow size, but all four other coarse axes also show more domain dependence
than expected from finite support alone. Inclusive thresholds retain boundary
ties. Consequently, tail sizes expand for discrete descriptors; in the most
extreme case the two disjoint DOCKSTRING ring-count tails jointly contain all
15,000 rows. Descriptor contrasts with an identical pair of tail sizes reuse
the same cached series of 200 random controls; the 14 table rows therefore do
not represent 14 independent control simulations.

\begin{table}[htbp]
\centering
\caption{Post-hoc descriptor-domain contrasts and matched-size random controls.
Control entries are medians [2.5th, 97.5th percentiles] across 200 disjoint
random-row splits with the same two tail sizes. They are finite-support
comparators, not confidence intervals or causal tests.}
\label{tab:s3descriptors}
\small
\resizebox{\textwidth}{!}{%
\begin{tabular}{llrrrl}
\toprule
Dataset & Descriptor & Low $n$ & High $n$ & Observed $\rho$ & Matched random control \\
\midrule
Docking-44 & Heavy atoms & 3,545 & 3,523 & 0.207 & 0.993 [0.988, 0.995] \\
 & Molecular weight & 3,163 & 3,163 & 0.205 & 0.992 [0.987, 0.994] \\
 & Labute ASA & 3,163 & 3,163 & 0.200 & 0.992 [0.987, 0.994] \\
 & TPSA & 3,163 & 3,167 & 0.753 & 0.992 [0.987, 0.994] \\
 & cLogP & 3,164 & 3,163 & 0.887 & 0.992 [0.988, 0.994] \\
 & Rotatable bonds & 3,736 & 3,168 & 0.789 & 0.993 [0.989, 0.995] \\
 & Ring count & 4,537 & 4,349 & 0.703 & 0.994 [0.992, 0.996] \\
\addlinespace
DOCKSTRING-58 & Heavy atoms & 4,367 & 3,817 & 0.361 & 0.991 [0.989, 0.993] \\
 & Molecular weight & 3,751 & 3,750 & 0.351 & 0.990 [0.988, 0.992] \\
 & Labute ASA & 3,750 & 3,750 & 0.322 & 0.990 [0.988, 0.992] \\
 & TPSA & 3,755 & 3,750 & 0.787 & 0.991 [0.988, 0.992] \\
 & cLogP & 3,750 & 3,750 & 0.822 & 0.990 [0.988, 0.992] \\
 & Rotatable bonds & 6,154 & 4,352 & 0.831 & 0.993 [0.991, 0.994] \\
 & Ring count & 8,376 & 6,624 & 0.803 & 0.995 [0.994, 0.996] \\
\bottomrule
\end{tabular}
}
\end{table}

Exact tail supports, thresholds and map statistics are in
\path{results/public_descriptor_domain_specificity/descriptor_extremes.csv};
family summaries and descriptor correlations are in
\path{results/public_descriptor_domain_specificity/descriptor_family_summary.csv}
and
\path{results/public_descriptor_domain_specificity/descriptor_correlations.csv}.
Matched-control replicates and summaries are in
\path{results/public_descriptor_domain_specificity/descriptor_random_disjoint_controls.csv}
and
\path{results/public_descriptor_domain_specificity/descriptor_random_disjoint_summary.csv}.

\section*{S4. Pilot recovery and panel-size sensitivity}

\subsection*{S4.1 Recovery of the source-library map}

For each dataset and calibration size, 100 simple-random pilot samples were
drawn without replacement. A pilot residual map was compared with the map on
all source rows not used by that pilot. Thus calibration and reference rows are
disjoint in every replicate. A 200-ligand pilot uses 1.58\% of Docking-44 score
cells and 0.077\% of DOCKSTRING score cells; the corresponding 500-ligand
fractions are 3.95\% and 0.192\%.

\begin{table}[htbp]
\centering
\caption{Pilot recovery on row-disjoint complement maps. Entries are replicate
means with central 95\% repeated-sample ranges over 100 pilots. Top-edge
columns report precision for the ten most positive or ten most negative
reference edges.}
\label{tab:s4edges}
\small
\resizebox{\textwidth}{!}{%
\begin{tabular}{llrrrr}
\toprule
Dataset & $n$ & Map $\rho$ & Sign agreement & Top positive 10 & Top negative 10 \\
\midrule
Docking-44 & 200 & 0.940 [0.888, 0.965] & 0.896 [0.850, 0.929] &
0.550 [0.200, 0.852] & 0.480 [0.200, 0.800] \\
Docking-44 & 500 & 0.973 [0.956, 0.982] & 0.931 [0.904, 0.951] &
0.618 [0.348, 0.900] & 0.606 [0.400, 0.800] \\
DOCKSTRING-58 & 200 & 0.920 [0.896, 0.936] & 0.877 [0.856, 0.895] &
0.639 [0.400, 0.852] & 0.677 [0.348, 0.900] \\
DOCKSTRING-58 & 500 & 0.965 [0.955, 0.972] & 0.920 [0.906, 0.936] &
0.779 [0.600, 0.900] & 0.776 [0.500, 0.952] \\
\bottomrule
\end{tabular}
}
\end{table}

Global map recovery therefore does not imply exact recovery of every extreme
edge. Full replicate summaries, including edge RMSE, top-10\% sign agreement
and top-25 lists, are in
\path{results/revision_map_decision_sensitivities/recovery_edge_summary.csv};
replicate-level edge results are in
\path{results/revision_map_decision_sensitivities/recovery_edge_metrics.csv}.
Including each calibration sample in the reference map changes the mean
agreements only to 0.941 and 0.975 for the 200- and 500-ligand Docking-44
pilots, and to 0.920 and 0.966 for DOCKSTRING; the corresponding row-disjoint
means in Table~\ref{tab:s4edges} are 0.940, 0.973, 0.920 and 0.965.

\subsection*{S4.2 Chemical-group-held-out recovery}

A stricter sensitivity formed five chemical-group folds and compared each
held-out fold's residual map with calibration maps sampled only from the other
four folds. Both datasets used one reproducible grouping rule after RDKit
cleanup, largest-fragment selection, uncharging, isotope removal and removal of
stereochemistry: canonical non-isomeric Bemis--Murcko scaffolds for cyclic
compounds and Standard-InChIKey connectivity blocks for acyclic compounds
(canonical non-isomeric SMILES fallback). This produced 5,112 groups among all
12,651 Docking-44 rows and 11,903 groups on the fixed seed-71 15,000-row
DOCKSTRING support. For
Docking-44, missing target scores were filled in every split using target means
estimated only on the four-fold calibration pool and then applied to both
calibration and held-out rows.

For the chemical-group design, 25 samples per held-out fold give 125 samples
per dataset and calibration size, with zero calibration--reference group
overlap throughout. A matched random-row design used the same held-out and
calibration sizes in each replicate. For Docking-44, that comparator
independently re-estimated target means on its random calibration pool before
filling either side, so held-out scores never informed imputation. Table~\ref{tab:s4group} shows that chemical-group holdout retained
high recovery, while all four paired group-minus-random central ranges included
zero. This tests group-held-out recovery within each source collection; it does
not establish transport to a shifted or external library.

For every held-out ligand, the exact maximum Tanimoto similarity to the full
calibration pool was computed from achiral radius-2, 2,048-bit Morgan
fingerprints on the same standardized parents. Median maxima across folds were
0.485--0.542 versus 0.630--0.644 for group and random Docking-44 holdouts, and
0.467--0.472 versus 0.491--0.493 for DOCKSTRING. Thus the group arm was
chemically harder by this diagnostic while retaining map recovery. The split
is not claimed to be fingerprint-cluster-disjoint; finite-bit Tanimoto equal to
one is not treated as proof of molecular identity.

\begin{table}[htbp]
\centering
\caption{Chemical-group-held-out map recovery and matched random-row holdout.
Entries are means [2.5th, 97.5th percentiles] across 125 repeated splits; the
paired column summarizes group minus matched-random recovery. These are
composition ranges, not confidence intervals or equivalence tests.}
\label{tab:s4group}
\small
\resizebox{\textwidth}{!}{%
\begin{tabular}{lrrll}
\toprule
Dataset & Calibration $n$ & Group-held-out $\rho$ & Matched-random $\rho$ & Paired difference \\
\midrule
Docking-44 & 200 & 0.925 [0.884, 0.956] & 0.935 [0.884, 0.960] & $-0.010$ [$-0.057$, 0.041] \\
Docking-44 & 500 & 0.958 [0.926, 0.975] & 0.968 [0.943, 0.980] & $-0.011$ [$-0.050$, 0.019] \\
DOCKSTRING-58 & 200 & 0.914 [0.889, 0.933] & 0.917 [0.893, 0.934] & $-0.003$ [$-0.034$, 0.029] \\
DOCKSTRING-58 & 500 & 0.958 [0.945, 0.965] & 0.960 [0.950, 0.968] & $-0.002$ [$-0.018$, 0.011] \\
\bottomrule
\end{tabular}
}
\end{table}

Summary and replicate rows are in
\path{results/public_scaffold_holdout_recovery/scaffold_holdout_summary.csv}
and
\path{results/public_scaffold_holdout_recovery/scaffold_holdout_metrics.csv};
fold diagnostics and exact nearest-neighbour records are in
\path{results/public_scaffold_holdout_recovery/chemical_group_fold_diagnostics.csv},
\path{results/public_scaffold_holdout_recovery/heldout_maximum_morgan_tanimoto.csv}
and
\path{results/public_scaffold_holdout_recovery/heldout_maximum_morgan_tanimoto_summary.csv};
the matched-random and paired summaries are in
\path{results/public_scaffold_holdout_recovery/matched_random_row_holdout_summary.csv}
and
\path{results/public_scaffold_holdout_recovery/chemical_group_minus_random_summary.csv}.

\subsection*{S4.3 A fixed panel-compression decision and sensitivity to $k$}

The decision example treats either positive or negative linear dependence as
computational redundancy and uses distance $d_{jk}=1-|r_{jk}|$. Pilot panels
are selected by deterministic BUILD+SWAP $k$-medoids; the fixed full-source
comparator is the exact mixed-integer $k$-medoids optimum. Both panels are
evaluated on each row-disjoint complement map. Consequently, the reported
difference is relative to one fixed full-source panel, not regret relative to a
complement-specific optimum, and can include BUILD+SWAP approximation error as
well as calibration-sample error. Average-linkage clustering on signed distance
$1-r$ is reported separately through adjusted Rand index (ARI).

\begin{table}[htbp]
\centering
\caption{Eight-target decision recovery. Values are means over 100 pilots;
brackets are central 95\% repeated-sample ranges. ``Random no worse'' is the
fraction of random eight-target panels with mean complement-map coverage
distance no greater than the pilot panel.}
\label{tab:s4decision}
\small
\resizebox{\textwidth}{!}{%
\begin{tabular}{llrrrr}
\toprule
Dataset & $n$ & ARI & Medoid overlap / 8 & Excess mean distance & Random no worse \\
\midrule
Docking-44 & 200 & 0.694 & 3.41 & 0.0153 [0.0037, 0.0325] & 0.0126 \\
Docking-44 & 500 & 0.784 & 4.63 & 0.0083 [0.0007, 0.0216] & 0.0006 \\
DOCKSTRING-58 & 200 & 0.595 & 3.65 & 0.0113 [0.0035, 0.0237] & 0.0010 \\
DOCKSTRING-58 & 500 & 0.658 & 5.01 & 0.0056 [0.0009, 0.0123] & 0.0001 \\
\bottomrule
\end{tabular}
}
\end{table}

The coverage result is not specific to $k=8$. Across $k=6,8,10$, mean excess
distance ranges from 0.01525 to 0.01604 at $n=200$ and from 0.00828 to 0.00868
at $n=500$ in Docking-44. DOCKSTRING ranges are 0.00994--0.01134 and
0.00488--0.00557. ARI ranges across those $k$ values are 0.690--0.712 and
0.745--0.785 for Docking-44 and 0.595--0.633 and 0.658--0.746 for
DOCKSTRING. These are ranges across three decision settings, not uncertainty
intervals. All $k$-specific means and repeated-sample ranges are in
\path{results/revision_map_decision_sensitivities/recovery_decision_summary.csv};
selected targets for every replicate are in
\path{results/revision_map_decision_sensitivities/recovery_medoid_selections.csv}.
The exercise quantifies compression of a Vina map and does not establish
biological target coverage or experimental selectivity.

\subsection*{S4.4 Raw and residual score-surface reconstruction}

Five chemical-group-disjoint folds tested whether selected raw target scores
reconstructed omitted target columns. Each fold used a 500-ligand pilot from
the other four folds; all scaling, exact residual-map $k$-medoids selection and
Ridge-GCV fitting were pilot-only. Table~\ref{tab:s4reconstruction} reports the
primary omitted-target variance-weighted $R^2$. The fold range is shown to make
the composition sensitivity visible; it is not a target-population confidence
interval.

\begin{table}[htbp]
\centering
\caption{Omitted-target reconstruction from residual-map-selected panels.
Entries are fold medians [minimum, maximum] across five chemical-group-disjoint
folds.}
\label{tab:s4reconstruction}
\small
\begin{tabular}{lrrr}
\toprule
Dataset & $k$ & Raw score surface & Row-centred residual surface \\
\midrule
Docking-44 & 4  & 0.881 [0.687, 0.886] & 0.234 [0.116, 0.249] \\
Docking-44 & 6  & 0.896 [0.737, 0.897] & 0.254 [0.145, 0.262] \\
Docking-44 & 8  & 0.900 [0.753, 0.910] & 0.278 [0.170, 0.283] \\
Docking-44 & 10 & 0.906 [0.762, 0.912] & 0.299 [0.178, 0.301] \\
Docking-44 & 12 & 0.907 [0.767, 0.913] & 0.316 [0.191, 0.319] \\
DOCKSTRING-58 & 4  & 0.650 [0.590, 0.675] & 0.184 [0.174, 0.192] \\
DOCKSTRING-58 & 6  & 0.682 [0.650, 0.707] & 0.202 [0.191, 0.207] \\
DOCKSTRING-58 & 8  & 0.698 [0.668, 0.726] & 0.219 [0.203, 0.221] \\
DOCKSTRING-58 & 10 & 0.705 [0.688, 0.731] & 0.228 [0.220, 0.231] \\
DOCKSTRING-58 & 12 & 0.711 [0.698, 0.743] & 0.242 [0.239, 0.246] \\
\bottomrule
\end{tabular}
\end{table}

Because different selectors omit different targets, their omitted-only scores
are not a common-outcome comparison. Table~\ref{tab:s4selectorcontrols}
therefore compares the corresponding all-target profiles, with observed panel
scores copied into their selected columns. The residual-map medoid selector was
competitive with pivoted QR on the raw surface and was better on the residual
surface in both datasets; the 50-panel random control did not remove the large
raw-versus-residual reconstruction gap.

\begin{table}[htbp]
\centering
\caption{Compact selector controls at $k=8$. Entries are median all-target
variance-weighted $R^2$ on five chemical-group-disjoint evaluation folds. The
two deterministic selectors contribute one value per fold; the fixed-random
row is the median over 50 panels per fold (250 panel--fold evaluations). These
all-target metrics provide a common outcome set across selectors; they are not
the selector-specific omitted-only values in Table~\ref{tab:s4reconstruction}.}
\label{tab:s4selectorcontrols}
\small
\begin{tabular}{llrr}
\toprule
Dataset & Pilot-only selector & Raw surface & Residual surface \\
\midrule
Docking-44 & Residual-map exact medoids & 0.918 & 0.391 \\
Docking-44 & Raw-correlation pivoted QR & 0.927 & 0.365 \\
Docking-44 & 50 fixed random panels & 0.904 & 0.359 \\
\addlinespace
DOCKSTRING-58 & Residual-map exact medoids & 0.740 & 0.315 \\
DOCKSTRING-58 & Raw-correlation pivoted QR & 0.716 & 0.290 \\
DOCKSTRING-58 & 50 fixed random panels & 0.716 & 0.280 \\
\bottomrule
\end{tabular}
\end{table}

At $k=8$, direct residual-outcome Ridge gave 0.278 and 0.219, essentially
identical to the residuals derived after raw-score prediction (differences
$1.5\times10^{-5}$ and $-6.3\times10^{-7}$). Selected-target PC1 baselines were
0.180 and 0.065. On common omitted targets, designed-minus-median-random mean
$R^2$ was positive in every fold for both matrices; medians were 0.036 and
0.058. Full profile, pooled RMSE, target-correlation, ligand-profile,
descriptor-only, pivoted-QR and 50-random-panel results are serialized in
\path{results/public_panel_selector_controls/}; the direct and common-omitted
controls are in \path{results/public_panel_selector_robustness/}.

\section*{S5. Experimental-map reliability and exact-support boundary}

\subsection*{S5.1 Repeatability of the DAVIS, PKIS2 and HotSpot maps}

All three experimental maps use the same 21 kinases and 210 target pairs. Two-way
residual maps were recomputed in 2,000 random disjoint half splits. The direct
half--half statistic estimates repeatability at the observed half-panel size;
it is not extrapolated to the full panel with Spearman--Brown. Whole-cluster
splits keep every Butina cluster in one half. Table~\ref{tab:s5reliability}
reports the empirical correlation estimator; Oracle Approximating Shrinkage
(OAS) is retained only as an estimator sensitivity \cite{chen2010}. These
splits retain all 72 DAVIS and 645 PKIS2 profiles, including the five and one
profiles that are constant across the 21 targets; the informative-profile
filter applies only to the cross-panel maps below. The Anastassiadis workbook
contains 178 compounds; two incomplete profiles were excluded, leaving 176.
Its primary split uses ligand rows because the workbook does not provide
structures. A separate PubChem crosswalk resolves 155 records for a secondary
cluster-weight sensitivity, with the remaining 21 treated as singletons.

\begin{table}[htbp]
\centering
\caption{Disjoint-half experimental residual-map repeatability. Brackets are
central 95\% repeated-split ranges over 2,000 splits, conditional on the fixed
21-target panel.}
\label{tab:s5reliability}
\small
\begin{tabular}{lrrr}
\toprule
Panel & Released ligands & Ligand-disjoint median [range] &
Butina-disjoint median [range] \\
\midrule
DAVIS & 72 & 0.757 [0.642, 0.840] & 0.753 [0.629, 0.840] \\
PKIS2 & 645 & 0.915 [0.874, 0.940] & 0.868 [0.807, 0.910] \\
HotSpot & 176 & 0.794 [0.690, 0.854] & Not primary \\
\bottomrule
\end{tabular}
\end{table}

Across ligand-disjoint splits, median edge-sign agreement is 0.767 in DAVIS
and 0.867 in PKIS2; median precision for the reference top correlation decile
is 0.619 and 0.714. The corresponding target-label null medians are close to
0.5 for sign agreement and 0.095 for top-decile precision. Split-level records
are in \path{results/experimental_map_reliability/map_reliability.csv};
edge-level bootstrap stability is in
\path{results/experimental_map_reliability/edge_stability.csv}. Bootstrap maps
overlap the full map, so those rows describe perturbation stability rather than
independent-map reliability.

The complete residual DAVIS and PKIS2 maps agree at Spearman $\rho=0.8484$
across 210 edges, with 79.52\% sign agreement and 57.14\% precision for the top
correlation decile. A joint target-label QAP with 9,999 permutations gives a
two-sided Monte Carlo probability of $1\times10^{-4}$, the resolution floor of
the design \cite{hubert1976}; null medians are $-0.0111$, 51.90\% and 9.52\%
for the three statistics. This is an external reproducibility
benchmark across different compounds and assay technologies, not an
attenuation-based noise ceiling. The complete row is
\path{results/experimental_map_reliability/cross_panel_geometry.csv}.

\subsection*{S5.2 Exact same-ligand support}

Full Standard InChIKey matching to DOCKSTRING yields 59 DAVIS and 154 PKIS2
ligands after duplicate cells are collapsed by their median. Three matched
DAVIS profiles are constant at the released pKd=5 floor, leaving 56 informative
primary matches; all 154 PKIS2 matches are informative. The broad comparator is
the 259,806-row DOCKSTRING map remaining after all 254 rows whose connectivity
blocks occur in DAVIS or PKIS2 are excluded.

\begin{table}[htbp]
\centering
\caption{Docking--experiment residual-map agreement on the primary informative
exact matches and on the chemically de-leaked broad DOCKSTRING reference.
Bootstrap brackets are paired Butina-cluster 95\% perturbation ranges for the
same-minus-broad difference, not target-population confidence intervals.}
\label{tab:s5same}
\small
\resizebox{\textwidth}{!}{%
\begin{tabular}{lrrrr}
\toprule
Panel & Exact informative ligands & Same support $\rho$ & Broad support $\rho$ &
Same minus broad [perturbation range] \\
\midrule
DAVIS & 56 & 0.3066 & 0.2995 & 0.0070 [$-0.1814$, 0.0904] \\
PKIS2 & 154 & 0.3261 & 0.3290 & $-0.0029$ [$-0.1338$, 0.0662] \\
\bottomrule
\end{tabular}
}
\end{table}

OAS correlations on the same supports are 0.3190 and 0.3223. In a further
whole-cluster disjoint-half check, the central 95\% repeated-split range for
same-half minus cross-half docking--experiment agreement is
$-0.0780$--0.0871 in DAVIS and $-0.0471$--0.0457 in PKIS2. Sharing one finite
ligand draw therefore does not explain a consistent gain over cross-half
comparison. Point estimates are in
\path{results/experimental_map_reliability/same_support_geometry.csv}, paired
bootstrap rows in
\path{results/experimental_map_reliability/same_support_bootstrap.csv}, and the
cross-half check in
\path{results/experimental_map_reliability/same_support_split_half.csv}.

These comparisons remove one chemical-support mismatch but do not validate
experimental affinity calibration or a biological co-selectivity network. No
classical correction for attenuation is applied because docking and experiment
cannot be assumed to be interchangeable measurements of one latent quantity.

\subsection*{S5.3 Cross-assay agreement after row centring}

Row centring increased edge-rank agreement for every assay pair: DAVIS--PKIS2
from 0.5821 to 0.8484, DAVIS--HotSpot from 0.6897 to 0.7651, and
PKIS2--HotSpot from 0.5264 to 0.7515. The comparison is paired over the same 210
target-pair edges. It shows greater reproducibility of a target-relative map
across these observed assays, not recovery of a true protein network or
independence among the three panels. HotSpot source validation, mappings,
excluded compounds, split-half rows and map-concordance records are in
\path{results/public_anastassiadis_panel_validation/}.

\subsection*{S5.4 Assay-value-blind fixed-target panel transfer}

Raw- and residual-map panels were selected from a 259,641-row DOCKSTRING
reference after excluding every DAVIS, PKIS2 and conservative HotSpot
connectivity. Every $\binom{21}{k}$ subset was evaluated for
$k=4,6,8,10,12$. Table~\ref{tab:s5transfer} reports the absolute signed
$1-r$ coverage losses of the lexicographic representative of each optimal
Vina or sequence panel and of the in-sample experimental oracle. It also
retains the conservative contrast, best raw optimum minus worst residual
optimum, across numerical Vina ties; positive values favour the residual
selector. Table~\ref{tab:s5panels} gives the representative panel compositions.

\begin{table}[htbp]
\centering
\caption{Absolute assay-value-blind signed co-response coverage losses on the
fixed 21-kinase panel. Lower absolute loss is better. Raw, residual and sequence
columns use the lexicographic representative among numerical optima; the oracle
uses the experimental map in sample and is a reference floor, not an
assay-value-blind selector. The final column is the conservative best-raw minus
worst-residual contrast across numerical Vina ties. All values are deterministic
for the stated maps; no assay value entered Vina panel selection.}
\label{tab:s5transfer}
\small
\begin{tabular}{lrrrrrr}
\toprule
Panel & $k$ & Raw Vina & Residual Vina & Sequence & Oracle &
Conservative contrast \\
\midrule
DAVIS & 4 & 0.6126 & 0.4799 & 0.4831 & 0.4240 & 0.1327 \\
DAVIS & 6 & 0.4437 & 0.4102 & 0.4014 & 0.3297 & 0.0235 \\
DAVIS & 8 & 0.4056 & 0.2657 & 0.3039 & 0.2508 & 0.1313 \\
DAVIS & 10 & 0.2932 & 0.2374 & 0.2515 & 0.1745 & 0.0471 \\
DAVIS & 12 & 0.2134 & 0.2025 & 0.1911 & 0.1154 & 0.0005 \\
\addlinespace
PKIS2 & 4 & 0.6519 & 0.6132 & 0.5704 & 0.5654 & 0.0387 \\
PKIS2 & 6 & 0.5589 & 0.5245 & 0.4956 & 0.4705 & 0.0344 \\
PKIS2 & 8 & 0.5057 & 0.4122 & 0.4145 & 0.3852 & 0.0935 \\
PKIS2 & 10 & 0.3975 & 0.3484 & 0.3381 & 0.3066 & 0.0491 \\
PKIS2 & 12 & 0.2887 & 0.2901 & 0.2717 & 0.2332 & $-0.0014$ \\
\addlinespace
HotSpot & 4 & 0.7000 & 0.5268 & 0.5773 & 0.5186 & 0.1543 \\
HotSpot & 6 & 0.5507 & 0.4452 & 0.4948 & 0.4284 & 0.0548 \\
HotSpot & 8 & 0.4787 & 0.3646 & 0.4171 & 0.3456 & 0.0936 \\
HotSpot & 10 & 0.3641 & 0.2980 & 0.3271 & 0.2689 & 0.0397 \\
HotSpot & 12 & 0.2691 & 0.2478 & 0.2517 & 0.1994 & 0.0186 \\
\bottomrule
\end{tabular}
\end{table}

\begin{table}[htbp]
\centering
\caption{Lexicographic representative panels for the primary signed coverage
objective. Numerical ties are resolved only for reproducible display; the
conservative inference in Table~\ref{tab:s5transfer} ranges over all tied Vina
optima. The same Vina and sequence panels are evaluated against each of the
three experimental maps.}
\label{tab:s5panels}
\scriptsize
\setlength{\tabcolsep}{3pt}
\begin{tabular}{rp{0.285\textwidth}p{0.285\textwidth}p{0.285\textwidth}}
\toprule
$k$ & Raw-Vina panel & Residual-Vina panel & Sequence-identity panel \\
\midrule
4 & CDK2, FGFR1, KDR, ROCK1 &
AKT1, JAK2, KDR, LCK &
AKT1, CSF1R, LCK, MAPK1 \\
6 & ABL1, CDK2, FGFR1, KDR, MET, ROCK1 &
AKT1, CSF1R, EGFR, JAK2, KDR, LCK &
AKT1, CSF1R, LCK, MAP2K1, MAPK14, MAPKAPK2 \\
8 & ABL1, CDK2, CSF1R, FGFR1, KDR, KIT, MET, ROCK1 &
AKT1, CSF1R, EGFR, JAK2, KDR, LCK, MAPKAPK2, PTK2 &
AKT1, CSF1R, EGFR, MAP2K1, MAPK14, MAPKAPK2, PLK1, SRC \\
10 & ABL1, CDK2, CSF1R, FGFR1, IGF1R, KDR, KIT, MAPK14, MET, ROCK1 &
AKT1, CSF1R, EGFR, IGF1R, JAK2, KDR, LCK, MAPK1, MAPKAPK2, PTK2 &
AKT1, CDK2, CSF1R, JAK2, MAP2K1, MAPK1, MAPKAPK2, MET, PLK1, SRC \\
12 & ABL1, AKT1, CDK2, CSF1R, FGFR1, IGF1R, KDR, KIT, LCK, MAPK1, MAPK14, MET &
AKT1, AKT2, CSF1R, EGFR, IGF1R, JAK2, KDR, LCK, MAPK1, MAPK14, MAPKAPK2, PTK2 &
AKT1, CDK2, CSF1R, EGFR, IGF1R, JAK2, MAP2K1, MAPK1, MAPKAPK2, PLK1, ROCK1, SRC \\
\bottomrule
\end{tabular}
\end{table}

The conservative contrast was positive in 14 of 15 cells. Normalized
trapezoidal AUCs over $k$ were 0.0671, 0.0489 and 0.0686 for DAVIS, PKIS2 and
HotSpot; their mean was 0.0616. A common 20,000-permutation target-label test
gave exploratory one-sided $p=0.0069$ after the reported
max-over-three-objectives adjustment. This post-hoc conditional randomization
result applies to the fixed 21-target set and does not provide confirmatory
error control for the full analysis history. In the HotSpot cluster multiplier,
all 2,000 perturbations were
positive (median 0.0628; 2.5--97.5\% range 0.0319--0.1063). Absolute and squared
distance objectives, CDK2 cyclin-A/cyclin-E alternatives, clipping,
target-median imputation, column-rank transformation, tie handling and the
locked two-panel versus triple-de-leaked reference are provided as sensitivity
records in the same result directory. After independently ranking ligands
within every experimental target before row centring, the primary signed
contrast remained positive in 14/15 cells (mean conservative contrast 0.0384;
max-over-three-objectives adjusted $p=0.0056$). Sequence was better than residual Vina
in 6 of 15 cells, so the comparison does not establish Vina superiority.

In a target-influence sensitivity, each of the 21 kinases was removed in turn
and both Vina selectors were re-optimized by complete enumeration on the
remaining 20-target maps. The joint normalized AUC contrast remained positive
for all 21 omissions (range 0.0294--0.0622), and 11--15 of the 15 assay--$k$
cells were positive for each omission. For DAVIS--PKIS2, DAVIS--HotSpot and
PKIS2--HotSpot, respectively, leave-one-target-out raw-map agreement ranged
0.541--0.699, 0.661--0.749 and 0.467--0.612; residual-map agreement ranged
0.798--0.875, 0.708--0.832 and 0.687--0.792. Residual agreement exceeded raw
agreement for every target omission in all three assay pairs. Complete records
are in
\path{results/public_anastassiadis_panel_validation/target_leave_one_out_panel_transfer.csv}
and
\path{results/public_anastassiadis_panel_validation/target_leave_one_out_map_concordance.csv}.
The 21 omissions overlap extensively and are reported as influence
sensitivities, not independent replications or a target-population interval;
single-target deletion also does not remove an entire redundant kinase group.

\section*{S6. Command-line audit and reproducibility}

\subsection*{S6.1 Public-core reconstruction}

As locked in \path{requirements.lock}, the public-core environment uses CPython
3.12.11, NumPy 1.26.4, pandas 2.3.2, SciPy 1.16.2, scikit-learn 1.7.2 and RDKit
2025.03.6 \cite{numpy2020,scipy2020,rdkit2025}. Environment specifications are
also in \path{requirements.txt} and \path{environment.yml}; \path{Dockerfile}
supplies a containerized alternative.
The following commands regenerate the analysis bundles and strict public-core
evidence ledger from the redistributed sources and checksum-fetched HotSpot
workbook:

\begin{verbatim}
.venv/bin/python analysis/cross_panel_overlap_audit.py

.venv/bin/python analysis/public_docking44_missing_data_sensitivity.py \
  --random-controls 1000 \
  --output-dir results/public_docking44_missing_data_sensitivity

.venv/bin/python analysis/public_residual_null_audit.py \
  --sample-size 12000 --repeats 500 \
  --output-dir results/public_residual_null_audit

.venv/bin/python analysis/public_chemical_domain_controls.py \
  --repetitions 200 \
  --output-dir results/public_chemical_domain_controls

.venv/bin/python analysis/public_descriptor_domain_specificity.py

.venv/bin/python analysis/revision_map_decision_sensitivities.py \
  --repetitions 100 --sizes 200,500 \
  --mw-null-repetitions 100 --cluster-k 6,8,10 \
  --output-dir results/revision_map_decision_sensitivities

.venv/bin/python analysis/public_scaffold_holdout_recovery.py

.venv/bin/python analysis/public_panel_selector_controls.py --overwrite

.venv/bin/python analysis/public_panel_selector_robustness.py --overwrite

.venv/bin/python analysis/experimental_map_reliability.py \
  --split-repeats 2000 --bootstraps 2000 --seed 20260810 \
  --output results/experimental_map_reliability

.venv/bin/python analysis/fetch_anastassiadis_supplement.py \
  data/source_cache/anastassiadis2011_moesm23.xls

.venv/bin/python analysis/public_anastassiadis_identity_audit.py --quiet

.venv/bin/python analysis/public_anastassiadis_panel_validation.py \
  --source data/source_cache/anastassiadis2011_moesm23.xls

.venv/bin/python analysis/build_public_core_evidence.py

.venv/bin/python analysis/build_public_core_macros.py

.venv/bin/python analysis/make_public_manuscript_figures.py \
  --output figures/public_core
\end{verbatim}

Arguments not shown use the producer defaults recorded in each bundle's
\texttt{summary.json} or \texttt{metadata.json}; changing a seed does not
reproduce the published replicate rows. The figure producer reads only the
strict-public ledger and cited bundles and renders with Matplotlib
\cite{matplotlib2007}.

Every multi-file result directory contains its own
\texttt{output\_checksums.json} manifest when
the producer emits multiple files. The evidence builder records every file it
opens in \path{results/public_core_source_manifest.csv}; this keeps the v4
report layer separate from historical ledgers. The focused verification suite
is:

\begin{verbatim}
.venv/bin/python -m pytest -q \
  analysis/test_cross_panel_overlap_audit.py \
  analysis/test_public_docking44_missing_data_sensitivity.py \
  analysis/test_public_residual_null_audit.py \
  analysis/test_public_chemical_domain_controls.py \
  analysis/test_public_descriptor_domain_specificity.py \
  analysis/test_revision_map_decision_sensitivities.py \
  analysis/test_public_scaffold_holdout_recovery.py \
  analysis/test_public_panel_selector_controls.py \
  analysis/test_public_panel_selector_robustness.py \
  analysis/test_experimental_map_reliability.py \
  analysis/test_public_anastassiadis_identity_audit.py \
  analysis/test_public_anastassiadis_panel_validation.py \
  analysis/test_dense_davis_benchmark.py \
  analysis/test_dense_pkis2_benchmark.py \
  analysis/test_residual_mechanism_analysis.py \
  analysis/test_target_map_audit.py \
  analysis/test_build_public_core_evidence.py \
  analysis/test_build_public_core_macros.py \
  analysis/test_v4_submission_consistency.py \
  analysis/test_make_public_manuscript_figures.py
\end{verbatim}

\subsection*{S6.2 Reusable target-map audit}

\path{analysis/target_map_audit.py} is a standalone CLI for a new dense
ligand-by-target numeric matrix. It reads no manuscript evidence ledger and
imports no repository-local analysis module. Target columns can be selected by
an explicit ordered list or regular expression. Missing values fail closed by
default and require an explicit target-mean or target-median imputation option;
malformed nonnumeric values always stop the run. A minimal invocation is:

\begin{verbatim}
python analysis/target_map_audit.py scores.csv \
  --output-dir results/my_target_map_audit \
  --target-regex '^score_' --ligand-id compound_id
\end{verbatim}

Domain comparison, row-disjoint pilot recovery and an eight-target compression
example can be requested together:

\begin{verbatim}
python analysis/target_map_audit.py scores.csv.gz \
  --output-dir results/my_target_map_audit \
  --target-regex '^target_' \
  --domain-column molecular_weight --domain-quantiles 4 \
  --pilot-sizes 100,200,500 --pilot-repeats 200 \
  --top-edge-count 10 --panel-k 8 --seed 20260810
\end{verbatim}

Each run writes a preprocessing contract, raw and residual target-correlation
matrices, the full eigenspectra and scalar summaries, reconstructable signed
edges, and a SHA-256 output manifest. Optional domain, pilot and medoid files
are emitted only when requested. The output directory must be absent or empty,
which prevents accidental mixing of artifacts from different contracts.
Complete CLI documentation is in \path{analysis/TARGET_MAP_AUDIT.md}; the
synthetic invariant tests verify that row centring preserves within-row target
ranks and that a controlled ligand-support shift changes the target map.

\begin{thebibliography}{99}
\small

\bibitem{vina2010} Trott O, Olson AJ (2010) AutoDock Vina: improving the speed and
accuracy of docking with a new scoring function, efficient optimization, and
multithreading. J Comput Chem 31:455--461. doi:10.1002/jcc.21334.

\bibitem{ding2023} Ding J, Tang S, Mei Z et al (2023) Vina-GPU 2.0: further
accelerating AutoDock Vina and its derivatives with graphics processing units.
J Chem Inf Model 63:1982--1998. doi:10.1021/acs.jcim.2c01504.

\bibitem{dockstring2022} Garc\'ia-Orteg\'on M, Simm GNC, Tripp AJ,
Hern\'andez-Lobato JM, Bender A, Bacallado S (2022) DOCKSTRING: easy molecular
docking yields better benchmarks for ligand design. J Chem Inf Model
62:3486--3502. doi:10.1021/acs.jcim.1c01334.

\bibitem{nikitin2025} Nikitin I, Morgunov I, Safronov V, Kalyuzhnaya A,
Fedorov M (2025) Towards explainable computational toxicology: linking
antitargets to rodent acute toxicity. Pharmaceutics 17:1573.
doi:10.3390/pharmaceutics17121573.

\bibitem{pan2003} Pan Y, Huang N, Cho S, MacKerell AD Jr (2003) Consideration of
molecular weight during compound selection in virtual target-based database
screening. J Chem Inf Comput Sci 43:267--272. doi:10.1021/ci020055f.

\bibitem{vigers2004} Vigers GPA, Rizzi JP (2004) Multiple active site corrections
for docking and virtual screening. J Med Chem 47:80--89. doi:10.1021/jm030161o.

\bibitem{jacobsson2006} Jacobsson M, Karl\'en A (2006) Ligand bias of scoring
functions in structure-based virtual screening. J Chem Inf Model 46:1334--1343.
doi:10.1021/ci050407t.

\bibitem{carta2007} Carta G, Knox AJS, Lloyd DG (2007) Unbiasing scoring
functions: a new normalization and rescoring strategy. J Chem Inf Model
47:1564--1571. doi:10.1021/ci600471m.

\bibitem{sepresa2009} Yang L, Luo H, Chen J, Xing Q, He L (2009) SePreSA: a
server for the prediction of populations susceptible to serious adverse drug
reactions implementing the methodology of a chemical--protein interactome.
Nucleic Acids Res 37:W406--W412. doi:10.1093/nar/gkp312.

\bibitem{casey2009} Casey FP, Pihan E, Shields DC (2009) Discovery of small
molecule inhibitors of protein--protein interactions using combined ligand and
target score normalization. J Chem Inf Model 49:2708--2717.
doi:10.1021/ci900294x.

\bibitem{luo2017} Luo Q, Zhao L, Hu J, Jin H, Liu Z, Zhang L (2017) The scoring
bias in reverse docking and the score normalization strategy to improve success
rate of target fishing. PLoS One 12:e0171433.
doi:10.1371/journal.pone.0171433.

\bibitem{lapillo2019} Lapillo M, Tuccinardi T, Martinelli A, Macchia M,
Giordano A, Poli G (2019) Extensive reliability evaluation of docking-based
target-fishing strategies. Int J Mol Sci 20:1023.
doi:10.3390/ijms20051023.

\bibitem{lee2020crds} Lee A, Kim D (2020) CRDS: consensus reverse docking system
for target fishing. Bioinformatics 36:959--960.
doi:10.1093/bioinformatics/btz656.

\bibitem{krause2023} Krause F, Voigt K, Di Ventura B, {\"O}zt{\"u}rk MA (2023)
ReverseDock: a web server for blind docking of a single ligand to multiple
protein targets using AutoDock Vina. Front Mol Biosci 10:1243970.
doi:10.3389/fmolb.2023.1243970.

\bibitem{fukunishi2005} Fukunishi Y, Mikami Y, Nakamura H (2005) Similarities
among receptor pockets and among compounds: analysis and application to
in silico ligand screening. J Mol Graph Model 24:34--45.
doi:10.1016/j.jmgm.2005.04.004.

\bibitem{fukunishi2006} Fukunishi Y, Mikami Y, Takedomi K, Yamanouchi M,
Shima H, Nakamura H (2006) Classification of chemical compounds by
protein--compound docking for use in designing a focused library. J Med Chem
49:523--533. doi:10.1021/jm050480a.

\bibitem{fukunishi2018} Fukunishi Y, Yamashita Y, Mashimo T, Nakamura H (2018)
Prediction of protein--compound binding energies from known activity data:
docking-score-based method and its applications. Mol Inform 37:e1700120.
doi:10.1002/minf.201700120.

\bibitem{lee2012profiles} Lee M, Kim D (2012) Large-scale reverse docking
profiles and their applications. BMC Bioinformatics 13(Suppl 17):S6.
doi:10.1186/1471-2105-13-S17-S6.

\bibitem{martin2011pqsar} Martin E, Mukherjee P, Sullivan D, Jansen J (2011)
Profile-QSAR: a novel meta-QSAR method that combines activities across the
kinase family to accurately predict affinity, selectivity, and cellular
activity. J Chem Inf Model 51:1942--1956. doi:10.1021/ci1005004.

\bibitem{bembenek2018} Bembenek SD, Hirst G, Mirzadegan T (2018)
Determination of a focused mini kinase panel for early identification of
selective kinase inhibitors. J Chem Inf Model 58:1434--1440.
doi:10.1021/acs.jcim.8b00222.

\bibitem{zhang2019informer} Zhang H, Ericksen SS, Lee CP et al (2019)
Predicting kinase inhibitors using bioactivity matrix derived informer sets.
PLoS Comput Biol 15:e1006813. doi:10.1371/journal.pcbi.1006813.

\bibitem{laufkotter2020} Laufk\"otter O, Laufer S, Bajorath J (2020)
Identifying representative kinases for inhibitor evaluation via systematic
analysis of compound-based target relationships. Eur J Med Chem 204:112641.
doi:10.1016/j.ejmech.2020.112641.

\bibitem{mangione2019cando} Mangione W, Samudrala R (2019) Identifying protein
features responsible for improved drug repurposing accuracies using the CANDO
platform: implications for drug design. Molecules 24:167.
doi:10.3390/molecules24010167.

\bibitem{kangas2014active} Kangas JD, Naik AW, Murphy RF (2014) Efficient
discovery of responses of proteins to compounds using active learning. BMC
Bioinformatics 15:143. doi:10.1186/1471-2105-15-143.

\bibitem{davis2011} Davis MI, Hunt JP, Herrgard S et al (2011) Comprehensive
analysis of kinase inhibitor selectivity. Nat Biotechnol 29:1046--1051.
doi:10.1038/nbt.1990.

\bibitem{wu2025davis} Wu MH (2025) A Complete and Modification-Aware DAVIS
Dataset. Harvard Dataverse, V3. doi:10.7910/DVN/RTQGP1.

\bibitem{drewry2017} Drewry DH, Wells CI, Andrews DM et al (2017) Progress
towards a public chemogenomic set for protein kinases and a call for
contributions. PLoS One 12:e0181585. doi:10.1371/journal.pone.0181585.

\bibitem{anastassiadis2011} Anastassiadis T, Deacon SW, Devarajan K,
Ma H, Peterson JR (2011) Comprehensive assay of kinase catalytic activity
reveals features of kinase inhibitor selectivity. Nat Biotechnol 29:1039--1045.
doi:10.1038/nbt.2017.

\bibitem{mazzucato2016} Mazzucato L, Fontanini A, La Camera G (2016) Stimuli
reduce the dimensionality of cortical activity. Front Syst Neurosci 10:11.
doi:10.3389/fnsys.2016.00011.

\bibitem{roy2007} Roy O, Vetterli M (2007) The effective rank: a measure of
effective dimensionality. In: 15th European Signal Processing Conference,
Pozna\'n, pp 606--610. doi:10.5281/zenodo.40328.

\bibitem{hubert1976} Hubert L, Schultz J (1976) Quadratic assignment as a
general data analysis strategy. Br J Math Stat Psychol 29:190--241.
doi:10.1111/j.2044-8317.1976.tb00714.x.

\bibitem{chen2010} Chen Y, Wiesel A, Eldar YC, Hero AO (2010) Shrinkage
algorithms for MMSE covariance estimation. IEEE Trans Signal Process
58:5016--5029. doi:10.1109/TSP.2010.2053029.

\bibitem{murcko1996} Bemis GW, Murcko MA (1996) The properties of known drugs.
1. Molecular frameworks. J Med Chem 39:2887--2893.
doi:10.1021/jm9602928.

\bibitem{morgan1965} Morgan HL (1965) The generation of a unique machine
description for chemical structures---a technique developed at Chemical
Abstracts Service. J Chem Doc 5:107--113. doi:10.1021/c160017a018.

\bibitem{rogers2010} Rogers D, Hahn M (2010) Extended-connectivity fingerprints.
J Chem Inf Model 50:742--754. doi:10.1021/ci100050t.

\bibitem{butina1999} Butina D (1999) Unsupervised data base clustering based on
Daylight's fingerprint and Tanimoto similarity: a fast and automated way to
cluster small and large data sets. J Chem Inf Comput Sci 39:747--750.
doi:10.1021/ci9803381.

\bibitem{rdkit2025} Landrum G et al (2025) RDKit: open-source cheminformatics
software, release 2025.03.6. Zenodo. doi:10.5281/zenodo.16996017.

\bibitem{numpy2020} Harris CR, Millman KJ, van der Walt SJ et al (2020) Array
programming with NumPy. Nature 585:357--362.
doi:10.1038/s41586-020-2649-2.

\bibitem{scipy2020} Virtanen P, Gommers R, Oliphant TE et al (2020) SciPy 1.0:
fundamental algorithms for scientific computing in Python. Nat Methods
17:261--272. doi:10.1038/s41592-019-0686-2.

\bibitem{matplotlib2007} Hunter JD (2007) Matplotlib: a 2D graphics environment.
Comput Sci Eng 9:90--95. doi:10.1109/MCSE.2007.55.

\end{thebibliography}


\end{document}
